% ============================================================================
% kapitel/iwt.tex
% ============================================================================

\section{Die grundlegende Größe: Information}

Die IWT hat nur eine fundamentale Größe: die Information \( I_k^{(n)} \).

\[
I_k^{(n)} \in \mathbb{C}
\]

\begin{itemize}
    \item \( k \) = Index des Informationsknotens
    \item \( n \) = diskreter Zeitindex
    \item \( I_k^{(n)} \) = komplexer Informationswert
\end{itemize}

Die Amplitude \( |I| \) kodiert die Stärke der Information.

Die Phase \( \arg(I) \) kodiert die kohärenten Relationen.

\section{Die Informationserhaltung}

Die Gesamtinformation ist erhalten:

\[
\sum_{k=1}^{N} |I_k^{(n)}|^2 = \text{const}
\]

Das ist Axiom 2 der IWT.

\section{Die Kopplungsmatrix}

Die Kopplungsmatrix beschreibt die Wechselwirkung zwischen Knoten:

\[
K_{kl} = \frac{1}{d_{kl}^{3-D}}
\]

Dabei ist \( d_{kl} \) die fraktale Distanz und \( D \) die fraktale Dimension.

\section{Die Dynamik}

Die Dynamik folgt aus dem Variationsprinzip:

\[
I_k^{(n+1)} = I_k^{(n)} + T \cdot \Phi_k
\]

Dabei ist \( \Phi_k \) die Kraft, die aus dem Lagrange-Funktional folgt.

\section{Die Rolle von Q}

Q ist nicht Teil der Axiome.

Q ist ein Resultat der Selbstorganisation.

Q emergiert aus der Information.